%------------------------------------------------------------------------------%
\section{Propriedades}

Nessa seção apresentamos algumas propriedades das séries que podemos
utilizar ao fazermos manipulações algébricas.

%------------------------------------------------------------------------------%
\begin{teorema}{Linearidade das Séries}{teo:linearidade-series}
  Se\index{Série!propriedades}
  \(\sum a_n\) e \(\sum b_n\)
  são \emph{séries convergentes} e $k\in\R$, então
  \begin{enumerate}[topsep=5mm,itemsep=5mm]
    \item \( \sum_{n=1}^\infty (a_n+b_n) =  \sum_{n=1}^\infty a_n +\sum_{n=1}^\infty b_n \),
    \item \( \sum_{n=1}^\infty (a_n-b_n) =  \sum_{n=1}^\infty a_n -\sum_{n=1}^\infty b_n \),
    \item \( \sum_{n=1}^\infty ka_n      = k\sum_{n=1}^\infty a_n \).
  \end{enumerate}
\end{teorema}
%------------------------------------------------------------------------------%

%------------------------------------------------------------------------------%
\begin{example}
  Calcular a soma da série
  \(\quad
    \sum_{n=0}^\infty \left(\frac{1}{2^n}+\frac{(-1)^n}{5^n}\right)
  \).
  \solution
  Observe que a série que queremos calcular é a soma, termo a termo de duas séries geométricas.
  Além disso, as razões das séries geométricas são $\abs{\sfrac{1}{2}}<1 $ e $\abs{\sfrac{-1}{5}}<1$.
  Então, podemos usar a regra da soma de séries e a fórmula da soma da série geométrica.
  \begin{align*}
    \sum_{n=0}^\infty \left(\frac{1}{2^n} +\frac{(-1)^n}{5^n}\right)
      & = \sum_{n=0}^\infty \left(\frac{ 1}{2}\right)^n
        + \sum_{n=0}^\infty \left(\frac{-1}{5}\right)^n \\[3mm]
      & = \frac{1}{1-\sfrac{1}{2}} + \frac{1}{1+\sfrac{1}{5}}
        = \frac{17}{6}   \,.
  \end{align*}
\end{example}
%------------------------------------------------------------------------------%

Note que, nesse exemplo, o índice começou em zero e não em 1. O início da série não tem
influência na convergência, mas pode mudar o valor da soma da série.

%------------------------------------------------------------------------------%
\begin{example}
  Calcular a soma da série
  \(\quad
    \sum_{n=1}^\infty\left[ \frac{1}{\sqrt{n}}
                          - \frac{1}{\sqrt{n+1}}
                          - \frac{1}{4^{n+1}}\right]
  \).
  \solution
  Observe que a série
  \(
    \sum_{n=1}^\infty\frac{1}{4^{n+1}}
  \)
  converge,
  pois é uma série geométrica com $r=\sfrac{1}{4}<1$, além disso,
  \[
    \sum_{n=1}^\infty\left(\frac{1}{\sqrt{n}}-\frac{1}{\sqrt{n+1}}\right)
  \]
  é uma série telescópica com a sequência de somas parciais dada por
  \[
    S_n = 1-\frac{1}{\sqrt{n+1}} \to 1  \,.
  \]
  Temos então que
  \begin{align*}
  	S & = \sum_{n=1}^\infty \left[\frac{1}{\sqrt{n}}-\frac{1}{\sqrt{n+1}}-\frac{1}{4^{n+1}}\right] \\
  	  & = \sum_{n=1}^\infty\left(\frac{1}{\sqrt{n}}-\frac{1}{\sqrt{n+1}}\right)
  	    - \sum_{n=1}^\infty\frac{1}{4^{n+1}}                            \\
  	  & = 1-\frac{1}{16}\sum_{n=1}^\infty\left(\frac{1}{4}\right)^{n-1} \\
  	  & = 1-\frac{1}{16}\;\frac{1}{1-\sfrac{1}{4}}                      \\
  	  & = 1 - \frac{1}{12} = \frac{11}{12}                              \,.
  \end{align*}
\end{example}
%------------------------------------------------------------------------------%

A hipótese de que $\sum a_n$ e $\sum b_n$ são séries convergentes é importante,
por exemplo, se considerarmos as duas séries divergentes
\[
  \sum_{n=1}^\infty (-1)^n
  \qquad \text{ e } \qquad
  \sum_{n=1}^\infty (-1)^{n+1}
\]
observamos que
\[
  \sum_{n=1}^\infty\big[(-1)^n+(-1)^{n+1}\big] = 0
  \quad\neq\quad
  \left(\sum_{n=1}^\infty (-1)^n\right) +
  \left(\sum_{n=1}^\infty (-1)^{n+1}\right)    \,.
\]

%------------------------------------------------------------------------------%
%------------------------------------------------------------------------------%
\subsection*{Manipulando Índices e Termos}
%------------------------------------------------------------------------------%
%------------------------------------------------------------------------------%

Primeiro vamos destacar uma característica do cálculo do limite de uma série,
não importa o valor dos primeiros termos, $n<N$, pois a existência do limite só
depende do que acontece no ``final'' da série. Assim quando formos verificar
a convergência sempre que for conveniente podemos ignorar o início da série.
Não se esqueça que isso só vale para determinar a convergência, o valor da
soma, se existir, provavelmente será alterado. O teorema a seguir formaliza
essa afirmação.

%------------------------------------------------------------------------------%
\begin{teorema}{Verificando a Convergência }{teo:verificando-convergencia}
  Ao verificarmos a convergência de uma série basta analisar a soma dos
  termos a partir de um índice $N$, isso é, se a série
  \[
    \sum_{n=N}^\infty a_n
  \]
  for convergente, ou divergente, então a série
  \[
    \sum_{n=0}^\infty a_n
  \]
  terá a mesma propriedade.
\end{teorema}
%------------------------------------------------------------------------------%
%\begin{proof}
%  Escrever a prova!
%\end{proof}
%------------------------------------------------------------------------------%

Podemos construir novas séries a partir de uma série dada, por exemplo,
adicionando ou retirando termos. Note que mudar a ordem dos termos de uma série
cria outra série. Se a mudança é feita em uma quantidade finita de termos
a nova série mantem a convergência ou divergência da série original.
Porém, se infinitos termos forem alterados será necessário
analisar separadamente a convergência da nova série.

Se a série
\[
  \sum_{n=0}^\infty a_n
\]
converge, então a série onde removemos, $k>0$, ou incluímos termos, $k<0$,
\[
  \sum_{n=k}^\infty a_n \qquad\qquad k\in \Z
\]
também converge, desde que os temos adicionais estejam bem definidos.
Se a série original divergir então a nova série também diverge.
Note que no caso de convergência, se $k\neq 0$ o limite é alterado.
De fato, se $L$ é o limite da série original,
no caso onde removemos termos, $k>0$, temos que
\[
  \sum_{n=k}^\infty a_n
  = \sum_{n=0}^\infty a_n - \left(a_1+\cdots+a_{k-1} \right)
  = L - \left(a_1+\cdots+a_{k-1}\right)                     \,.
\]
Enquanto que, quando acrescentando termos, $k<0$, temos que
\[
  \sum_{n=k}^\infty a_n
  = \sum_{n=0}^\infty a_n + \left( a_k+\cdots+a_{-1} \right)
  = L + \left( a_k+\cdots+a_{-1} \right)        \,.
\]

%------------------------------------------------------------------------------%
\begin{proposicao}{Alterando os Termos de uma Série}{prop:numero-termos-serie}
  Adicionar, retirar ou permutar um \emph{número finito de termos} em uma série,
  não altera o fato da série convergir ou divergir.

  Porém, pode alterar o valor da soma da série, caso ele exista.
\end{proposicao}
%------------------------------------------------------------------------------%

%------------------------------------------------------------------------------%
\begin{example}
  Verifique a convergência da série
  \(
    \sum_{n=-3}^\infty \frac{1}{2^n}
  \).
  \solution
  A série
  \[
    \sum_{n=0}^\infty \frac{1}{2^n}
  \]
  converge, pois é uma série geométrica com $r=\sfrac{1}{2}$.
  Portanto, a série
  \[
    \sum_{n=-3}^\infty \frac{1}{2^n}
  \]
  também converge. Além disso,
  \begin{align*}
    \sum_{n=-3}^\infty \frac{1}{2^n} & = 2^3 + 2^2 + 2^1 + \sum_{n=0}^\infty \frac{1}{2^n} \\[3mm]
                                     & = 8 + 4 + 2 + \frac{1}{1-\sfrac{1}{2}}              \\
                                     & = 14 + 2                                            \\
                                     & = 16                                                \,.
  \end{align*}
\end{example}
%------------------------------------------------------------------------------%

%------------------------------------------------------------------------------%
\begin{proposicao}{Reindexando uma Série}{prop:reindexando-serie}
  Alterar o índice inicial de uma série da seguinte maneira
  \[
    \sum_{n=0}^\infty a_n =
    \sum_{n=k}^\infty a_{n-k} \qquad \forall\; k \in \Z
  \]
  não altera a convergência ou a soma da série, caso ela exista.
\end{proposicao}
%------------------------------------------------------------------------------%

Note que, quando $k<0$, estamos iniciando a série com índices menores do que zero.
Por outro lado, se $k>0$, estamos começando a série com índices positivos.
Claro que esse tipo de mudança nos índices também podem ser feitas em séries
indexadas a partir de~$n\neq 0$. De forma geral temos
\[
  \sum_{n=m}^\infty a_n \;=\!\!\sum_{n=m+k}^\infty a_{n-k}               \,.
\]
Note que não é necessário decorar essa fórmula, basta fazer um mudança de variáveis.

%------------------------------------------------------------------------------%
\begin{example}
  Reescreva a série
  \(
    \sum_{n=1}^\infty\frac{5}{n(n+1)}
  \)
  indexando-a a partir de $-1$.
  \solution
  Como $n$ começa em $1$ e queremos indexar a partir de $-1$,
  basta introduzir um novo índice $m=n-2$, então $n=m+2$ e
  $m=-1, \dots, \infty$.
  Assim, a série indexada a partir de $-1$ é
  \[
    \sum_{n=1}^\infty\frac{5}{n(n+1)} =
    \sum_{m=-1}^\infty\frac{5}{(m+2)(m+3)}   \,.
  \]
\end{example}
%------------------------------------------------------------------------------%

%------------------------------------------------------------------------------%
\begin{example}
  Reindexe a série do exemplo anterior a partir $20$.
  \solution
  Para indexar a partir de $20$, queremos que
  $m=20$ corresponda a $n=1$, ou seja,
  quando $m=20$ devemos ter $m+a=n=1$ e $m+b=n+1=2$.
  Assim os valores de $a$ e $b$ devem ser $-19$ e $-18$.
  Concluímos que a nova série é escrita como
  \[
    \sum_{n=1}^ \infty\frac{5}{n(n+1)} =
    \sum_{m=20}^\infty\frac{5}{(m-19)(m-18)}   \,.
  \]
\end{example}
%------------------------------------------------------------------------------%

%------------------------------------------------------------------------------%
